\documentclass[11pt]{article}
\usepackage[utf8]{inputenc}
\usepackage[T1]{fontenc}
\usepackage[margin=1in]{geometry}
\usepackage{xcolor}
\usepackage{booktabs}
\usepackage[colorlinks=true,linkcolor=blue,urlcolor=blue]{hyperref}
\definecolor{cProblem}{RGB}{192,26,26}
\definecolor{cMethod}{RGB}{15,118,110}
\definecolor{cResult}{RGB}{30,122,60}
\definecolor{cGap}{RGB}{194,120,0}

\begin{document}

\section*{Research Paper Template: What Goes in Each Section}

\textit{Use this template to see what belongs in each part of a research paper. The colors show the job of each idea. The running example is ECG image digitization: turn a photo or scan of a paper ECG into a clean digital signal.}

\fbox{\textbf{How to read the colors}}

\begin{itemize}
\item \colorbox{cProblem}{\textcolor{white}{\textbf{ Problem }}} \textcolor{cProblem}{red} marks the problem you solve.
\item \colorbox{cMethod}{\textcolor{white}{\textbf{ Method }}} \textcolor{cMethod}{teal} marks the method or idea.
\item \colorbox{cResult}{\textcolor{white}{\textbf{ Result }}} \textcolor{cResult}{green} marks the result or claim.
\item \colorbox{cGap}{\textcolor{white}{\textbf{ Gap }}} \textcolor{cGap}{orange} marks the gap in past work.
\end{itemize}

\fbox{\textbf{The whole paper at a glance}}

\begin{tabular}{ll}
\toprule
\textbf{Section} & \textbf{Its one job} \\
\midrule
1. Title & Name the work in one clear line \\
2. Abstract & The whole paper in one short paragraph \\
3. Introduction & Walk from the big picture to your question \\
4. Related Work & Show past work and where you fit \\
5. Methodology & Explain your method so others can rebuild it \\
6. Experiments & Describe the test setup \\
7. Results & Report what the numbers say \\
8. Discussion & Explain what the results mean \\
9. Limitations & Say what your work does not cover \\
10. Conclusion & Wrap up and point forward \\
11. References & List every source you cited \\
\bottomrule
\end{tabular}

\section{Title}
\textit{Purpose: name the paper so a reader knows the topic in one line.}

\fbox{\textbf{What to put here}}
\begin{itemize}
\item The \textcolor{cMethod}{method or idea} and the \textcolor{cProblem}{problem} it solves.
\item Keep it short, plain, and specific. No hype words.
\end{itemize}

\fbox{\textbf{ECG example}}

\textcolor{cMethod}{Robust ECG Image Digitization} \textcolor{cProblem}{Under Common Image Damage}: A Method by Method Benchmark on PhysioNet 2024

\section{Abstract}
\textit{Purpose: give the whole paper in one short paragraph.}

\fbox{\textbf{What to put here}}
\begin{itemize}
\item One or two lines of context, then the \textcolor{cProblem}{problem}.
\item What you did (the \textcolor{cMethod}{method}) and the main \textcolor{cResult}{result}.
\item Aim for 4 to 6 short sentences. No citations, no tables.
\end{itemize}

\fbox{\textbf{ECG example}}

\textcolor{cProblem}{Paper ECG scans often carry creases, shadows, and ink smudges.} These flaws can lower the quality of the digital signal. \textcolor{cGap}{Past work reports one overall score and does not show how each type of damage changes the result.} \textcolor{cMethod}{We test three digitization methods on PhysioNet 2024 and group the images by damage type.} \textcolor{cResult}{We find that the best method changes with the damage type, and no single method wins on every case.}

\section{Introduction}
\textit{Purpose: walk the reader from the big picture down to your exact question.}

\fbox{\textbf{What to put here}} An introduction is built in a clear order. Follow these steps.

\begin{itemize}
\item \textbf{Step 1. Wide context.} Say why the area matters. \\
ECG: many clinics still keep heart records only on paper.
\item \textbf{Step 2. The problem.} Name the specific pain. \\
ECG: \textcolor{cProblem}{scanned paper ECGs are often damaged by creases, shadows, and smudges.}
\item \textbf{Step 3. What is missing.} Say what current work does not do. \\
ECG: current tools report a single overall accuracy score.
\item \textbf{Step 4. The gap.} Name the gap in one clear line. \\
ECG: \textcolor{cGap}{no study shows if accuracy holds up for each separate type of damage.}
\item \textbf{Step 5. The question.} Turn the gap into a question. \\
ECG: does digitization accuracy hold up for each type of image damage, method by method?
\item \textbf{Step 6. The contributions.} List what you add.
\begin{itemize}
\item \textcolor{cMethod}{We benchmark three digitization methods, one damage type at a time.}
\item \textcolor{cResult}{We report SNR in dB for each method and each damage type.}
\item \textcolor{cResult}{We show which method to trust for which kind of damage.}
\end{itemize}
\end{itemize}

\section{Related Work}
\textit{Purpose: show what others did and where your work sits.}

\fbox{\textbf{What to put here}}
\begin{itemize}
\item Group the prior \textcolor{cMethod}{methods} and say what each one does.
\item Give the limit of each group, then point to the \textcolor{cGap}{gap}.
\item Compare the works. Do not just list them one by one.
\end{itemize}

\fbox{\textbf{ECG example}}

\textcolor{cMethod}{Early methods trace the ECG line with hand written rules.} \textcolor{cMethod}{Newer methods use deep networks to read the signal.} Both report one overall score on clean or lightly damaged images. \textcolor{cGap}{None breaks the score down by damage type, which is the gap we fill.}

\section{Methodology}
\textit{Purpose: explain your method so another person could rebuild it.}

\fbox{\textbf{What to put here}}
\begin{itemize}
\item The \textcolor{cMethod}{method} step by step, with inputs and outputs.
\item The model or pipeline, and the settings you used.
\item Enough detail that a reader could repeat your work.
\end{itemize}

\fbox{\textbf{ECG example}}

\textcolor{cMethod}{Our pipeline has five steps: load the ECG image, remove the printed grid, find the signal line, map the pixels to millivolts, and save the digital signal.} We run this same pipeline for each of the three methods so the comparison is fair.

\section{Experiments}
\textit{Purpose: describe the setup you used to test the method.}

\fbox{\textbf{What to put here}}
\begin{itemize}
\item The dataset, the benchmark, and the \textcolor{cResult}{metric} you measure.
\item The baselines you compare against.
\item How you split the data. This is the plan, not the numbers yet.
\end{itemize}

\fbox{\textbf{ECG example}}

We use the PhysioNet 2024 benchmark. The metric is \textcolor{cResult}{SNR in dB}, where higher is better. We compare three methods, called A, B, and C. We split the images into five damage groups: clean scan, creases, shadow, ink smudge, and low resolution.

\section{Results}
\textit{Purpose: report what the numbers say, with no opinion yet.}

\fbox{\textbf{What to put here}}
\begin{itemize}
\item The \textcolor{cResult}{results} in clear tables, grouped the way you planned.
\item A direct answer to your question, backed by the numbers.
\end{itemize}

\fbox{\textbf{ECG example}} \textbf{Table: example SNR in dB, higher is better. Numbers are illustrative.}

\begin{tabular}{lccc}
\toprule
\textbf{Damage type} & \textbf{Method A} & \textbf{Method B} & \textbf{Method C} \\
\midrule
Clean scan & 24.8 & 24.1 & 22.0 \\
Creases & 18.2 & 20.4 & 15.1 \\
Shadow & 16.7 & 19.9 & 14.0 \\
Ink smudge & 12.1 & 17.6 & 10.5 \\
Low resolution & 21.3 & 20.0 & 19.2 \\
\bottomrule
\end{tabular}

\textcolor{cResult}{Method B gives the highest SNR on creases, shadow, and ink smudge. Method A wins on clean scans and low resolution.}

\section{Discussion}
\textit{Purpose: explain what the results mean and why.}

\fbox{\textbf{What to put here}}
\begin{itemize}
\item Read the \textcolor{cResult}{result} and say what it tells us.
\item Link back to the \textcolor{cGap}{gap} and your question.
\item Say when the method works well and when it does not.
\end{itemize}

\fbox{\textbf{ECG example}}

\textcolor{cResult}{The best method depends on the type of damage.} This \textcolor{cGap}{answers the gap directly}: one overall score hides these swings and can point a user to the wrong method for their damaged scans.

\section{Limitations}
\textit{Purpose: say honestly what your work does not cover.}

\fbox{\textbf{What to put here}}
\begin{itemize}
\item Where the \textcolor{cMethod}{method} breaks or was not tested.
\item Limits of the data and the setup. Be specific and honest.
\end{itemize}

\fbox{\textbf{ECG example}}

We test five damage types on one benchmark only. \textcolor{cProblem}{We do not test images that carry two kinds of damage at the same time,} so mixed damage is still an open case.

\section{Conclusion}
\textit{Purpose: wrap up in a few lines and point forward.}

\fbox{\textbf{What to put here}}
\begin{itemize}
\item Restate the \textcolor{cProblem}{problem} and the main \textcolor{cResult}{result} in short.
\item One line on useful future work.
\end{itemize}

\fbox{\textbf{ECG example}}

\textcolor{cProblem}{Damaged ECG scans are hard to digitize.} \textcolor{cResult}{We showed that the best method changes with the type of damage.} Future work should test mixed damage and add more methods.

\section{References}
\textit{Purpose: list every source you cited, in one steady style.}

\fbox{\textbf{What to put here}}
\begin{itemize}
\item A full citation for each work: authors, title, venue, and year.
\item Use the same style for all entries. List only sources you actually cited.
\item Use the url command for web links so they stay clickable.
\end{itemize}

\fbox{\textbf{ECG example}}
\begin{itemize}
\item PhysioNet 2024 Challenge. Digitization of ECG images. 2024. \url{https://physionet.org}
\item A. Author, B. Author. Deep ECG signal recovery from paper scans. Journal of Biomedical Signals, 2023.
\end{itemize}

\end{document}
